In this paper, we introduced Neural State-Dependent Delays Differential Equations (Neural SDDDE), an auto-differentiable framework for solving Delay Differential Equations (DDEs) with multiple and state-dependent delays via neural networks. This open-source, robust python DDE solver is written in JAX, hardware agnostic and pushes the current envelope of Neural DDEs by handling the delays in a more generic way. To the best of our knowledge, no machine learning library or open-sourced code are available to model in general fashion DDEs. 

The need for such a tool rises in modelling and prediction of dynamical systems describing real-life problems is apparent. Indeed, many systems are modelled leveraging the Markovianity property, but this is often only an approximation of more realistic descriptions where the effects of past states affect the current ones. So far, standard ways to data-drive dynamical models via neural networks are by circumventing these limitations with lifting the state to a higher dimensional space – for instance applying the augmented version of Neural ODE (ANODE) – or considering different paradigms such as Neural Laplace. From the theoretical viewpoint, with respect to ANODEs, Neural DDEs preserve physical interpretability of the state vector and allow the detection of the time delays. With respect to Neural Laplace, a technique encompassing a broad number of different DEs, Neural DDEs are capable of handling delays that are state-dependent and do not necessitate a large amount of data for the initialization of the training. Finally, as already mentioned, the current toolbox generalizes the current envelope of DDEs, where only a single constant delay or several piece-wise constant delays can be handled. 

In this work, NODEs, the augmented version ANODE and Neural Laplace are used for benchmarking Neural SDDDE using several numerical experiments of incremental complexity, including the Lokta-Volterra and Mackey-Glass models, time- and state- dependent models, and a system characterized by delay diffusion. The dynamics of these models is correctly reproduced for all these cases by Neural SDDDE. Moreover, we have shown that Neural SDDDE performs well compared these with established methods in terms of accuracy and reliability, and is highly effective in modeling DDEs, in particular state-dependent ones when given the delays.

In future research, we plan to expand upon this approach by studying an equivalent version of ANODE with Neural SDDDEs, along with the examination of DDEs of neutral type (NDDEs). Moreover, we believe this flexible and versatile tool may provide a valuable contribution to several fields such as control theory where time-delay are often considered. In particular, it may prove useful in learning a model for partially observed systems whose dynamics of observables can be learned, under mild conditions, from their time-history. 


% \\ \\
% \emph{We believe this flexible and versatile tool may provide a valuable contribution to several fields. In particular, it may prove useful in learning a model for partially observed systems whose dynamics of observables can be learned, under mild conditions, from their time-history via a memory-capable model. Prototypical of this framework is the Takens' embedding theorem representation, grounded in the Mori-Zwanzig formalism, which justifies an efficient functional structure for representing the dynamics of observables with delayed differential equations. We are actively exploring the use of the present tool for learning the governing equations, with relevant non-constant delays, of the measurements of a weakly-observed flow configuration.}